\begin{center}
  \textbf{\LARGE Summary Sheet}\\[0.35em]
  \textbf{Team Control Number: \# 2625838}
\end{center}

\setlength{\parskip}{-0.5em}
\setlength{\parindent}{0pt}
\setlist[itemize]{itemsep=0.10em, topsep=0.1em, parsep=0em, partopsep=0em}

\small
We study \emph{Dancing with the Stars} (DWTS) as an inverse problem: judges' scores and weekly eliminations are observed
across 34 seasons\cite{comap2026_problemC_data}, but audience vote totals are withheld. The key unknown is each active contestant's weekly fan vote share,
which drives eliminations under different aggregation rules.

\subsection*{What we built}
\begin{itemize}
    \item A reproducible season--week--contestant panel that standardizes active sets, elimination labels, and judge-score aggregates
    under irregular show formats (no-elimination weeks, multi-elimination weeks, varying numbers of judges, and post-exit ``zero'' encodings).
    \item A convex inference model (Task A) that reconstructs latent fan vote shares consistent with observed eliminations, and an identifiability
    measure (KPI2) quantifying how uniquely votes are determined each week.
    \item A replay engine (Tasks B/C) that re-simulates every season under percent-based versus rank-based aggregation, and faithfully implements the incumbent ``combined Bottom-2 + judges' save'' procedure as the baseline for controlled counterfactual comparisons.
    \item Mixed-effects attribution models (Task D1) separating drivers of judges' scores and fan support, and a professionalism-first institutional redesign (Task D2) via \textit{Professional Bottom-2} (\textit{PB2}):
    judges-only gating defines the Bottom-2 risk set (judge Bottom-2), and fan influence is applied only \emph{within} that gate via a transparent blended score; late-stage variants allow structured suspense (bottom-$k$ band / limited multi-elimination),
    validated by replay-based stress tests: (i) pairwise Jaccard similarity across Percent/Rank/B2/PB2, (ii) controversial-case replays, and (iii) a unified bias test.
\end{itemize}

\subsection*{Sanity-check evidence (data integrity)}
From the weekly panel we define \emph{eligible} season--weeks (including multi-elimination weeks evaluated by bottom-$k$ set match).
Under a BL-0 baseline (judge-only with a uniform fan proxy), the elimination hit-rate is 0.364 (rank aggregation) and 0.375 (percent aggregation)
across 264 eligible weeks, and the two schemes disagree on the predicted eliminated set in 3.8\% of eligible weeks (FlipRate $=0.038$).
These rule divergences are consistent with classical observations in voting and aggregation theory.\cite{saari2001decisions}
These metrics validate consistent labeling and replay mechanics and serve as a baseline---not the final inferred-vote results. When counterfactual replays require post-exit inputs, we use a conservative, minimal-information imputation
only to complete the replay, and we do not interpret these imputed segments as forecasting contestant performance.

\subsection*{Key deliverables for decision-making}
\begin{itemize}
    \item Season-by-season comparisons of how often rank vs.\ percent rules change eliminations and winners, and when judges' save overturns vote-driven outcomes.
    \item Uncertainty reports showing which weeks/contestants are non-identifiable from observed eliminations alone.
    \item A recommended voting policy with interpretable control knobs (judge-gated risk set, judge--fan balance within the gate, and late-stage suspense strength),
    selected/validated via replay-based stress tests (Jaccard similarity, controversial-case replays, and a unified bias test).
\end{itemize}

\subsection*{Management recommendation}
Pilot the proposed \textit{Professional Bottom-2} (PB2) rule in a one-season ``shadow mode'' (computed in parallel with the broadcast rule) and adopt an operating point supported by replay-based stress tests
on eligible weeks while improving professionalism, as evidenced by cross-rule Jaccard comparisons, controversial-case replays, and a lower bias score under PB2.
